\section{Related Work}
\label{sec:related}

\paragraph{Food classification and detection.}
Early work cast food understanding as whole-image classification.
Food-101~\cite{bossard2014food101} established the canonical benchmark with $101$ dishes and $1$k images per class and has driven deep-network baselines for over a decade.
UECFood-100 and UECFood-256~\cite{matsuda2012uecfood} added bounding-box annotations so that multiple items on a tray can be localized.
These datasets treat each dish as an atomic unit, which works well for restaurant menus but poorly for home-cooked plates where a single dish is itself a mixture (e.g., a stir-fry).
SmartPlate targets the mixed-plate regime explicitly and therefore operates at the level of \emph{ingredient}, not dish.

\paragraph{Food segmentation.}
Pixel-level food understanding started with small, single-cuisine datasets such as UNIMIB2016 and became tractable at scale with FoodSeg103~\cite{wu2021foodseg103}, which provides $7{,}118$ images spanning $103$ ingredient classes with dense masks.
Subsequent work has pushed quality with stronger backbones (DeepLabV3+~\cite{chen2018deeplabv3plus}) and promptable segmentation (FoodSAM~\cite{lan2023foodsam}, which composes the Segment Anything Model~\cite{kirillov2023sam} with a food classifier).
We fine-tune SegFormer~\cite{xie2021segformer} on FoodSeg103 because its hierarchical MiT encoder and lightweight MLP decoder strike a favorable accuracy/parameter trade-off and fit comfortably on a laptop GPU through the MPS backend.

\paragraph{Nutrition estimation.}
Prior pipelines for nutrition estimation fall into two camps.
Dataset-driven regression, exemplified by Nutrition5k~\cite{thames2021nutrition5k}, uses RGB-D recordings of dishes with ground-truth ingredient masses to directly regress calories and macros.
Rule-based pipelines, instead, recognize ingredients (via classification or segmentation) and look up a reference table~\cite{usdafdc}.
Regression can be more accurate when depth or turntable imagery is available but is brittle outside the training distribution and does not naturally expose \emph{why} a plate is judged the way it is.
SmartPlate follows the second line: we estimate ingredient proportions from a single RGB image and look up per-group macros from USDA FoodData Central~\cite{usdafdc}, keeping the final score a transparent function of observable quantities.

\paragraph{Dietary scoring rubrics.}
Operationalizing ``healthy eating'' as a score has a long history in nutrition epidemiology, most notably the Healthy Eating Index family used by the USDA.
We adopt the Harvard Healthy Eating Plate~\cite{harvardplate} because its visual target---half vegetables and fruits, a quarter whole grains, a quarter healthy protein---maps naturally onto the per-class pixel proportions produced by a semantic segmentation model.
To our knowledge this pairing of a public dietary rubric with a laptop-scale food segmenter has not been reported in prior computer-vision work.
